%% Chapter 2 — Literature Survey
%% Guideline target: ~2 pages.
%%
%% Keep this tight. In Phase 2 the literature survey is NOT the main deliverable
%% — reviewers have already seen your Phase 1 survey. Summarise only what is
%% directly relevant to the architecture and methodology you actually built.
%%
%% Style convention for this template: one subsection per paper, each ending
%% with a one-line "Relevance to our work" sentence.

\chapter{Literature Survey}
\label{ch:literature}

% ---------------------------------------------------------------------------
\section{Paper Title Goes Here}
\label{sec:lit-paperA}

\textbf{Authors:} Author et al. \quad \textbf{Year:} 20XX \quad
\textbf{Venue:} Journal / Conference

\textit{[Replace.] One paragraph summarising what this paper does, the
dataset used, and the headline result. Do not copy from the paper — paraphrase.}

\textit{[Replace.] One paragraph on limitations of the paper and how your
project addresses or extends them.}

% Example of an in-line figure reference (image must exist under figures/):
% \begin{figure}[htbp]
%   \centering
%   \includegraphics[width=0.8\linewidth]{figures/lit_paperA_arch.png}
%   \caption{Architecture proposed by Author et al.}
%   \label{fig:lit-paperA-arch}
% \end{figure}

% ---------------------------------------------------------------------------
\section{Second Paper Title}
\label{sec:lit-paperB}

\textit{[Replace with content, structured like the section above.]}

% ---------------------------------------------------------------------------
\section{Research Gaps}
\label{sec:lit-gaps}

\textit{[Replace.] Summarise the gaps identified across the surveyed work
that motivate your project's contribution. Bullet points are acceptable here.}

\begin{itemize}
  \item \textit{[Gap 1]}
  \item \textit{[Gap 2]}
  \item \textit{[Gap 3]}
\end{itemize}
